%% toc: P15 164 164.1,164.11
\chapter{Relevance Axioms, Their Provenance, and Carnap's Relevance Measure}
\label{sec:p15-w02b-p15-024-relevance}

Chapter~\ref{sec:p15-rationality-4-3} sets out Arthur Merin's account of
evidential relevance: the qualitative relations, the measures D, Q, L and
their transforms, the ordinal characterization of the three principal
measures, the task-relative ranking results, and the composition rules.
This chapter supplies five further parts of the same 2021 chapter and
of the works it depends on. The first is Merin's argument for the status
of the four axioms (p.~271), with the two inferences that argument uses,
proved in full. The second is his account of where the measures and
axioms come from, set against the different attributions in Crupi and
Tentori's 2016 survey and in Carnap's 1950 book. The third is a
comparison of Merin's probabilistic relevance with Carnap's relevance
theory in Chapter~\ref{sec:p15-carnap-confirmation-relevance-and-reliability}
and with Hempel's objection to the raven solution Merin adopts. The
fourth takes the positions Merin summarizes in a sentence, on old
evidence, common causes, diagnostic categories and measure pluralism,
and proves the probability facts that bound each of them. The fifth is
the status of each dependent result: the relative-distance
measure Z, whose characterization Merin cites, is characterized and
proved here, and every work the chapter depends on is listed with its
standing \cite[pp.~265--275]{Merin2021Relevance}.

The inputs are a finite probability model, or a regular measure in
Carnap's sense, and one or more source and target propositions. The
conclusions are of three kinds: ordinal facts about which measures
satisfy which axioms; identities linking Carnap's relevance measure to
Merin's covariance measure; and the dependence of the raven ordering on
a frequency assumption. Notation follows
Definitions~\ref{def:p15-r43-polarity}, \ref{def:p15-r43-extreme}
and~\ref{def:p15-r43-measures}: \(R_T(E\Yright H)_P\) is the relevance
of source \(E\) to target \(H\) under \(P\) on a measure of type \(T\),
and \(\Delta_P(A,B)=P(AB)-P(A)P(B)\).

\begin{hbconvention}{def:p15-w02b-p15-024-claim-kinds}{Four kinds of claim}
\begin{enumerate}
\item \emph{Formal}: statements about defined measures on finite
  probability models or regular measures. Each is stated as a numbered
  result and proved in the annex.
\item \emph{Normative or methodological}: recommendations about which
  measure or axiom a user should accept, such as Merin's preference for
  measures that satisfy POP-D. These are positions argued in the source.
  The formal results below show what accepting one axiom commits a user
  to; they do not by themselves settle which axiom to accept.
\item \emph{Historical}: who first defined or named a measure or axiom.
  These are reported as the attributions of a named writer, with that
  writer's locator. The primary works they cite are not reproduced here,
  and where two cited writers disagree both attributions are given.
\item \emph{Empirical}: claims about what people judge, what
  commentators accept, or how a practice came about. The chapter
  contains three, all reported by Merin without data: that most
  commentators treat POP-D as mandatory; that some probability
  judgments are covert relevance judgments
  (Section~\ref{sec:p15-w02b-p15-024-surveys}); and Pearl's thesis that
  diagnostic categories are formed to screen off symptoms
  (Section~\ref{sec:p15-w02b-p15-024-common-cause}). None is supported
  here by a study, population, design or uncertainty estimate.
\end{enumerate}
\end{hbconvention}

\section{Which comparisons each axiom governs}
\label{sec:p15-w02b-p15-024-axiom-roles}

The four axioms of Section~\ref{sec:p15-r43-axioms} fall into two
pairs according to what is held fixed. POP-D and MCIE compare two
sources \(E_1,E_2\) for one target \(H\). LIK-D and MDHD compare two
targets \(H_1,H_2\) for one source \(E\). Merin draws the consequence
that POP-D and LIK-D are not rivals: they legislate for different
comparisons, and the quotient measure Q satisfies both
\cite[p.~271]{Merin2021Relevance}. The characterization in
Theorem~\ref{thm:p15-r43-characterization} confirms that the pair
POP-D and LIK-D is consistent, since it singles out the increasing
transforms of Q.

Crupi and Tentori state the same four conditions with the hypothesis
written first, \(C_P(h,e)\) for the confirmation of \(h\) by \(e\)
under \(P\), and under different names. Merin's note~11 records the
correspondence: POP-D is their \emph{Final probability}, LIK-D their
\emph{Law of likelihood}, MCIE their \emph{Modularity (for
conditionally independent data)}, and MDHD their \emph{Disjunction of
alternative hypotheses}. The survey numbers them (C1) to (C4) in that
order and adds the formality assumption (F): the measure is a function
of \(P(h\land e)\), \(P(h)\) and \(P(e)\) alone
\cite[p.~274, n.~11]{Merin2021Relevance}
\cite[pp.~654--655]{CrupiTentori2016Confirmation}. Formality is the
part of Merin's Transparency condition that the characterization
actually uses; smoothness is not needed
\cite[pp.~269, 271]{Merin2021Relevance}.

\subsection{Merin's case for the status of the axioms}
\label{sec:p15-w02b-p15-024-axiom-case}

Merin ranks the axioms as follows \cite[p.~271 and p.~274,
nn.~12--14]{Merin2021Relevance}.

\begin{enumerate}
\item \emph{POP-D.} He reports that most commentators regard it as
  mandatory and that denying it feels counterintuitive. He adds that
  Steel (2003) argued POP-D to be necessary for enumerative induction
  that proceeds from the proportion of favourable to unfavourable
  outcomes in a series of structurally identical experiments. From
  these he concludes that only measures satisfying POP-D are attractive
  as measures of first resort. The first ground is a report about
  opinion, given without a survey. The second is an attribution: the
  chapter states Steel's conclusion but not his argument, so here it
  has the standing of a cited thesis, and no result below depends on it.
\item \emph{LIK-D.} Merin notes that LIK-D is often called the Law of
  Likelihood, and that philosophers, unlike statisticians, apply the
  likelihood concept also to hypotheses that are compatible and not
  exhaustive (note~14). He judges LIK-D less compelling than POP-D, but
  adds the caveat that intuitions about single axioms carry limited
  normative weight. The reason is that the axioms interact: a user who
  accepts POP-D and regards direction assignments as unwarranted is
  thereby committed to LIK-D, and a user who accepts POP-D but wants
  asymmetric values, at least for extreme relevance, must accept
  violations of LIK-D. Theorem~\ref{thm:p15-w02b-p15-024-axiom-status}
  proves both inferences.
\item \emph{MCIE.} Merin locates MCIE in Good (1960) and Heckerman
  (1988) and values it for its special case, the multiplicativity of L
  and additivity of LL for evidence that is conditionally independent
  given the target and given its negation (his Lemma~3, here
  Theorem~\ref{thm:p15-r43-composition}). Crupi and Tentori say that
  their label for this condition is freely adapted from Heckerman (1988,
  pp.~18--19), and that the name ``law of likelihood'' goes back to
  Hacking (1965) \cite[p.~655, n.~7]{CrupiTentori2016Confirmation}.
  Merin adds that these practical benefits are unavailable for other
  increasing transforms of L, naming KO. Section~\ref{sec:p15-w02b-p15-024-lemma3}
  makes the claim exact: only positive powers of L keep the product
  rule, only positive multiples of LL keep the sum rule, and KO keeps
  neither (Theorem~\ref{thm:p15-w02b-p15-024-ko-transform}).
\item \emph{MDHD.} It looks sensible but forces D to give a refuted
  hypothesis and its weakening by an irrelevant disjoint disjunct the
  same decrement; Section~\ref{sec:p15-r43-axioms} works the card
  example, which is also Crupi and Tentori's illustration
  \cite[p.~657]{CrupiTentori2016Confirmation}. Merin's remark that Q
  and L do not face the problem, so that D is ordinally inequivalent
  to each, is proved with the L values in
  Theorem~\ref{thm:p15-w02b-p15-024-mdhd-ratios}.
\end{enumerate}

\begin{hbtheorem}{thm:p15-w02b-p15-024-axiom-status}{proved=annex:p15-w02b-p15-024-axiom-status}{What POP-D commits a user to}
Let \(R\) be a relevance measure given by one formal function of the
source--target joint table, compared only on events whose probabilities
lie strictly between zero and one. Call \(R\) \emph{amount-symmetric}
if \(R(A\Yright B)_P=R(B\Yright A)_P\) for all such \(A,B,P\).
\begin{enumerate}
\item If \(R\) satisfies POP-D and is amount-symmetric, it satisfies
  LIK-D.
\item If \(R\) satisfies POP-D and LIK-D, it is amount-symmetric.
  Consequently a measure satisfying POP-D that assigns different values
  to \(A\Yright B\) and \(B\Yright A\) for one such pair, for instance
  where \(A\) is extremely positive to \(B\) but not conversely,
  violates LIK-D.
\item D satisfies POP-D and violates LIK-D. Let \(k\) be uniformly
  distributed on the thirty integers \(15,\ldots,44\), and let \(A\),
  \(B\) and \(N\) say that \(15\), \(3\) and \(9\) divide \(k\),
  respectively. Then \(P(A)=1/15\), \(P(B)=1/3\), \(P(N)=1/10\), and
  \(A\) and \(N\) each entail \(B\). Hence
  \(P(B\mid A)=P(B\mid N)=1\), while
  \[
   R_D(B\Yright A)_P=\tfrac{2}{15},\qquad
   R_D(B\Yright N)_P=\tfrac15,\qquad
   R_D(A\Yright B)_P=\tfrac23 .
  \]
  The quotient gives \(R_Q(B\Yright A)_P=R_Q(B\Yright N)_P=3\).
\end{enumerate}
\end{hbtheorem}

The model in item~3 is a finite realization of Merin's divisibility
example in Section~\ref{sec:p15-r43-rankings}, which leaves the
distribution over numbers open. The first two items show why Merin
discounts verdicts on single axioms: whether LIK-D is acceptable is
settled, given POP-D, by the separate question whether relevance should
have a direction. The same items explain the user advice in Merin's
conclusion, which recommends Q or QL where direction is intrinsically
unspecified \cite[p.~274]{Merin2021Relevance}.

\subsection{What Lemma~3 gives L, and what MDHD costs D}
\label{sec:p15-w02b-p15-024-lemma3}

Merin's Lemma~3 states that, for evidence conditionally independent
given \(H\) and given \(\neg H\), the likelihood ratio of the
conjunction is the product of the two likelihood ratios, and its
logarithm the sum. He then says that these practical benefits of MCIE
are unavailable for other increasing transforms of L, and names KO
\cite[p.~271]{Merin2021Relevance}. Two remarks fix the claim. First,
the printed sum in Lemma~3 carries the subscript of L on its first
summand; the identity that holds, and the one proved in
Theorem~\ref{thm:p15-r43-composition}, has LL in both summands.
Second, ``unavailable'' cannot mean that KO values of a conjunction are
undetermined by the KO values of the parts: they are determined, by a
rule that is neither a sum nor a product. Note~7 of the chapter writes
KO as a hyperbolic tangent of half a logarithm
\cite[p.~274, n.~7]{Merin2021Relevance}; the logarithm must be taken
of L, so that \(R_{KO}=\tanh(\tfrac12 R_{LL})\), which is the form in
Theorem~\ref{thm:p15-r43-measures}. Writing KO this way is what
exposes its composition rule.

\begin{hbtheorem}{thm:p15-w02b-p15-024-ko-transform}{proved=annex:p15-w02b-p15-024-ko-transform}{Which increasing transforms of L keep Lemma~3}
Let \(0<P(H)<1\), and for \(i=1,2\) let \(a_i=P(E_i\mid H)\) and
\(b_i=P(E_i\mid\neg H)\) lie strictly between zero and one, with
\((E_1\perp E_2\mid\pm H)_P\). Write \(l_i=R_L(E_i\Yright H)_P\) and
\(k_i=R_{KO}(E_i\Yright H)_P\).
\begin{enumerate}
\item \(k_i=(l_i-1)/(l_i+1)=\tanh\bigl(\tfrac12R_{LL}(E_i\Yright H)_P\bigr)\),
  \(R_L(E_1E_2\Yright H)_P=l_1l_2\), and
  \[
   R_{KO}(E_1E_2\Yright H)_P=\frac{k_1+k_2}{1+k_1k_2}.
  \]
\item In the model with \(P(H)=1/2\), \(a_1=a_2=3/4\),
  \(b_1=b_2=1/4\) and conditional independence given \(H\) and given
  \(\neg H\): \(l_1=l_2=3\) and \(R_L(E_1E_2\Yright H)_P=9\), while
  \(k_1=k_2=1/2\) and \(R_{KO}(E_1E_2\Yright H)_P=4/5\), which differs
  from both \(k_1+k_2=1\) and \(k_1k_2=1/4\).
\item Let \(f\colon(0,\infty)\to\mathbb{R}\) be strictly increasing and
  put \(R_T=f(R_L)\) on such models. Then
  \(R_T(E_1E_2\Yright H)_P=R_T(E_1\Yright H)_P+R_T(E_2\Yright H)_P\)
  holds in every such model if and only if \(f=c\ln\) for some
  \(c>0\), and
  \(R_T(E_1E_2\Yright H)_P=R_T(E_1\Yright H)_P\,R_T(E_2\Yright H)_P\)
  holds in every such model if and only if \(f(x)=x^c\) for some
  \(c>0\).
\end{enumerate}
\end{hbtheorem}

So Merin's claim is true in this exact sense: among increasing
transforms of L, the product rule survives only for positive powers of
L, which are L itself on a changed scale, and the sum rule only for
positive multiples of LL, which is LL in another logarithmic base. KO
is neither; its composition rule, item~1, is the addition law of the
hyperbolic tangent. The practical loss is that KO values of
independent reports cannot be pooled by adding or multiplying them; one
must first undo the transform.

Merin also says that the MDHD embarrassment of D does not arise for L
and Q, and that this shows D to be ordinally inequivalent to each
\cite[p.~271]{Merin2021Relevance}. Chapter~\ref{sec:p15-rationality-4-3}
gives the Q values for the card example. The L values and the general
case are as follows.

\begin{hbtheorem}{thm:p15-w02b-p15-024-mdhd-ratios}{proved=annex:p15-w02b-p15-024-mdhd-ratios}{The MDHD case on D, Q and L}
Let \(0<P(E)<1\), let \(E\,\mathrm{exneg}_P\,H\), and let \(G\) satisfy
\(P(G)>0\), \(P(HG)=0\), \((E\perp G)_P\) and \(P(H\vee G)<1\). Write
\(h=P(H)\) and \(g=P(G)\). Then:
\begin{enumerate}
\item \(R_D(E\Yright H)_P=R_D(E\Yright H\vee G)_P=-h\);
\item \(R_Q(E\Yright H)_P=0\) and \(R_Q(E\Yright H\vee G)_P=g/(h+g)>0\);
\item \(R_L(E\Yright H)_P=0\) and
  \[
   R_L(E\Yright H\vee G)_P=\frac{g\,(1-h-g)}{(h+g)(1-g)}>0;
  \]
\item \(E\) is negative but not extremely negative to \(H\vee G\).
\end{enumerate}
Consequently D ties the two targets while Q, L and every strictly
increasing transform of either (QL, LL, KO) ranks \(H\vee G\) strictly
above \(H\); D is ordinally inequivalent to each of them in the sense
of Definition~9 of the source. In the fifty-two-card model of
Chapter~\ref{sec:p15-rationality-4-3}, with \(E\) a face card, \(H\)
the seven of spades and \(G\) a red card, \(h=1/52\) and \(g=1/2\), so
the D values are both \(-1/52\), the Q values are \(0\) and \(26/27\),
and the L values are \(0\) and \(25/27\).
\end{hbtheorem}

Item~4 is the content of Merin's phrase that \(E\) speaks conclusively
against \(H\) but only inconclusively against the weakened hypothesis.
The L values also show what the weakening does: refuting \(H\) moves L
to its minimum, while the disjunction keeps an L value below one by an
amount fixed by \(h\) and \(g\) alone, and that amount shrinks to zero
as \(h\) does.

\subsection{Source rankings and the direction of measurement}
\label{sec:p15-w02b-p15-024-direction}

Merin's Lemma~4 says that if POP-D is a soundness requirement on
relevance measures, then a sound ordinal ranking of several sources
\(E_1,\ldots,E_n\) by their relevance to one target \(H\) does not
depend on the direction of measurement. He derives it from his
Theorem~9 (here the first paragraph of
Theorem~\ref{thm:p15-r43-rankings}) together with the amount symmetry
of Q \cite[p.~271]{Merin2021Relevance}. The lemma is stated in words;
the result below fixes what ``sound'' and ``direction'' have to mean
for it to be true, and shows where it stops.

\begin{hbtheorem}{thm:p15-w02b-p15-024-direction}{proved=annex:p15-w02b-p15-024-direction}{Merin's Lemma~4 in a precise form}
Fix a probability \(P\) on a finite set, a target \(H\) and sources
\(E_1,\ldots,E_n\), \(n\geq2\), all with probabilities strictly between
zero and one. Call a ranking of the sources \emph{sound} if it is the
ranking, ties included, given by \(R(E_i\Yright H)_P\) for some
measure \(R\) that is one formal function of the source--target joint
table, as in Theorem~\ref{thm:p15-w02b-p15-024-axiom-status}, and
satisfies POP-D.
\begin{enumerate}
\item There is exactly one sound ranking: \(E_i\) is ranked above,
  level with or below \(E_j\) exactly when \(P(H\mid E_i)\) is
  greater than, equal to or less than \(P(H\mid E_j)\).
\item For every \(i\),
  \(R_Q(E_i\Yright H)_P=R_Q(H\Yright E_i)_P\). Hence the sound ranking
  is also the ranking by the reversed quotients \(R_Q(H\Yright E_i)_P\),
  and by the reversed values \(R(H\Yright E_i)_P\) of any such measure
  \(R\) satisfying both POP-D and LIK-D.
\item The reversed values of a measure satisfying POP-D alone need not
  give the sound ranking. With \(P(H)=1/2\), \(P(E_1)=1/20\),
  \(P(H\mid E_1)=4/5\), \(P(E_2)=1/2\) and \(P(H\mid E_2)=7/10\),
  the sound ranking puts \(E_1\) above \(E_2\), while
  \(R_D(H\Yright E_1)_P=3/100<1/5=R_D(H\Yright E_2)_P\).
\end{enumerate}
\end{hbtheorem}

So the lemma holds in the sense that the ranking itself is fixed by the
target and the posteriors, and that a direction-free measure, Q or any
measure satisfying POP-D and LIK-D, computes it whichever way its
arguments are written. It does not say that every measure satisfying
POP-D can be read in either direction: item~3 is a counterexample with
D, the measure that satisfies POP-D and violates LIK-D
(Theorem~\ref{thm:p15-w02b-p15-024-axiom-status}). This is the
limitation that leads Merin to recommend Q or QL where direction is
unspecified \cite[p.~274]{Merin2021Relevance}. The lemma also concerns
only rankings of sources for one target; for one source and several
targets, the second part of Theorem~\ref{thm:p15-r43-rankings} shows
that D, Q and L disagree.

\section{Provenance of the measures, the notes, and two open questions}
\label{sec:p15-w02b-p15-024-provenance}

\subsection{Who defined and named the measures}
\label{sec:p15-w02b-p15-024-history}

Merin attaches the following history to the measures of
Definition~\ref{def:p15-r43-measures}
\cite[pp.~270, 272--273 and p.~274, n.~8]{Merin2021Relevance}.
\begin{itemize}
\item D is Hugh MacColl's \(\delta\), of 1897 or earlier, with a
  reference to Keynes's \emph{Treatise} of 1921, pp.~154--155.
\item Q is W.~E. Johnson's \emph{coefficient of influence}, dated 1921
  or earlier and published in Keynes's \emph{Treatise}. Its other name,
  \emph{relevance quotient}, is Carnap's of 1950 (note~8). Some authors
  call it the probability ratio measure and label it \(R\); Merin warns
  that this label already had other referents.
\item L, up to logarithms, was commended by C.~S. Peirce, A.~M. Turing
  and I.~J. Good. Merin bases its information-theoretic primacy on two
  facts: it is sufficient for the data's bearing on the bipartition of
  H and its negation, and every other sufficient statistic determines
  it (Savage 1954, pp.~128--140).
  Theorem~\ref{thm:p15-r43-sufficiency} proves the finite binary case.
\item The logarithmic quotient QL is from Good (1950); the sigmoid KO
  from Kemeny and Oppenheim (1952); Z from Crupi and Tentori (2013).
\item The covariance measure \(R_C\) and the definition of extreme
  positivity are credited to Carnap (1950); the definition is marked
  ``modified'' \cite[pp.~266, 272]{Merin2021Relevance}.
\item Eells and Fitelson (2002) are cited for the large number of
  measures in the literature and for the prevailing view that target
  negation should always invert relevance while source negation should
  not \cite[pp.~269, 272]{Merin2021Relevance}.
\item Good (1983) applied the term \emph{quasi-utility} to LL, with a
  different motivation from Merin's use of it for all three measures
  \cite[p.~273]{Merin2021Relevance}.
\end{itemize}

Two of the works cited here check parts of this. Carnap's relevance chapter
introduces the relevance quotient under that name, ascribes the concept
and its theorems to W.~E. Johnson and Keynes, and says that Keynes
called it the coefficient of influence on the basis of Johnson's
unpublished notes (Keynes, pp.~150--155). Carnap also names the
quotient and the difference of posterior and prior confirmation as the
two obvious ways to measure positive relevance, before choosing a third
function, his relevance measure \(r\)
\cite[\S66, pp.~356--357; \S67, p.~362]{Carnap1950Probability}.
Crupi and Tentori give a different priority account: the probability
difference first defined by Carnap, the probability ratio by Keynes,
and the likelihood ratio by Turing as Good (1950) reports it
\cite[p.~654, n.~5]{CrupiTentori2016Confirmation}. Their attribution
of the difference to Carnap and Merin's attribution of it to MacColl
cannot both describe the first definition. Neither of the two texts quotes
MacColl, so the question is left as a disagreement between the two
attributions.

\subsection{The notes that qualify the main text}
\label{sec:p15-w02b-p15-024-notes}

Note~1 says that much of the chapter could be carried out with other
representations of degrees of credence, and points to the chapters by
H\'ajek and Staffel, by Dubois and Prade, and by Kern-Isberner,
Skovgaard-Olsen and Spohn in the same handbook. These correspond to
Chapters~\ref{sec:p15-rationality-4-1}, \ref{sec:p15-rationality-4-7}
and~\ref{sec:p15-rationality-5-3}. The claim is qualified (``to a good
extent'') and no transfer is carried out in the chapter; every result
here and in Chapter~\ref{sec:p15-rationality-4-3} is stated for
probability only \cite[p.~265 and p.~274, n.~1]{Merin2021Relevance}.

Note~5 concerns language. Glosses on the likelihood \(P(D\mid H)\)
often describe it counterfactually, as the probability that \(D\)
would have had if \(H\) were true; Merin refers to Edwards (1972) for
likelihoods in statistics \cite[p.~274, n.~5]{Merin2021Relevance}.
The gloss and the credal conditional come apart exactly in the
old-evidence case of Section~\ref{sec:p15-r43-old-evidence}.

\begin{hbtheorem}{thm:p15-w02b-p15-024-certain-data}{proved=annex:p15-w02b-p15-024-certain-data}{Certain data have uninformative credal likelihoods}
If \(P(D)=1\), then \(P(D\mid H)=1\) for every \(H\) with
\(P(H)>0\). If moreover \(0<P(H)<1\), then
\(P(D\mid\neg H)=1\) and \(R_L(D\Yright H)_P=1\).
\end{hbtheorem}

So once \(D\) is known, the credal state assigns every hypothesis the
same likelihood. A likelihood that discriminates between hypotheses in
that state has to come from elsewhere, for instance from a statistical
model for hypothetical repetitions, which is the alternative Merin
attributes to working statisticians \cite[p.~268]{Merin2021Relevance}.
The counterfactual wording names that other quantity; it is not the
conditional probability of the current credal state.

Note~17 connects the fact that D gives the less probable of two mutually
relevant propositions the larger relevance to the other (stated in
Theorem~\ref{thm:p15-r43-measures}) with what Roberto Festa calls
\emph{Matthew properties}, citing Festa and Cevolani (2016). The note
gives only the biblical gloss of the name, not a definition, so the
connection is given here as Merin's pointer
\cite[p.~272 and p.~274, n.~17]{Merin2021Relevance}.

Note~20 names a second application of evidential relevance that, unlike
the reading of experimental instructions in Section~\ref{sec:p15-r43-use},
involves no misunderstanding of the question: the semantics of function
words such as \emph{and}, \emph{or}, \emph{not}, \emph{but},
\emph{even}, \emph{also}, \emph{some} and \emph{many}, developed in
Merin (1999). Note~16 gives the one fragment of that programme the
chapter states, the comparison of inclusive and exclusive \emph{or},
which is proved in Theorem~\ref{thm:p15-r43-composition}
\cite[p.~274, nn.~16, 20]{Merin2021Relevance}.

\subsection{The accounting analogy and the combined-report question}
\label{sec:p15-w02b-p15-024-open-questions}

Merin motivates the measures with three questions about rain
\cite[p.~268]{Merin2021Relevance}. Let \(A_1\) say there are dark
clouds, \(A_2\) that the forecast predicts rain, \(A_3\) that it
predicts dry weather, and \(B\) that it will rain. Which of \(A_1,A_2\)
speaks more strongly for \(B\)? Does \(A_1A_3\) speak for \(B\),
against it, or neither? Is rain or snow more strongly favoured by dark
clouds? He says that none of these is usually settled by qualitative
considerations, and that rankings may depend on the measure, much as a
ranking of companies by asset value may depend on the accounting
convention.

The analogy has a precise limit. Theorem~\ref{thm:p15-r43-rankings}
shows that the dependence is confined to some tasks: for one target and
several sources (the first question) every measure satisfying POP-D gives
the same ranking, while for one source and several targets (the third
question) D, Q and L can disagree. The second question is different
again, since it concerns the sign of a compound source, and the signs of
its parts do not fix it.

\begin{hbtheorem}{thm:p15-w02b-p15-024-combined-report}{proved=annex:p15-w02b-p15-024-combined-report}{Part polarities do not fix the polarity of the conjunction}
Let \(P(B)=1/2\), make \(A_1\) and \(A_3\) conditionally independent
given \(B\) and given \(\neg B\), and put \(P(A_1\mid B)=9/10\),
\(P(A_1\mid\neg B)=3/10\), \(P(A_3\mid\neg B)=1-P(A_3\mid B)\).
Then \(A_1\) is positive and \(A_3\) negative to \(B\) whenever
\(P(A_3\mid B)<1/2\), and
\[
\begin{array}{c|ccc}
P(A_3\mid B)&2/5&1/4&1/5\\\hline
R_L(A_3\Yright B)_P&2/3&1/3&1/4\\
P(B\mid A_1A_3)&2/3&1/2&3/7\\
\text{polarity of }A_1A_3\text{ to }B&\text{positive}&\text{irrelevant}&\text{negative}
\end{array}
\]
\end{hbtheorem}

The three models agree on the qualitative facts Merin's question
supplies and differ on the answer. What decides it here is the product
\(R_L(A_1\Yright B)_P\,R_L(A_3\Yright B)_P\), which is meaningful only
under the conditional-independence premise of
Theorem~\ref{thm:p15-r43-composition}.

\section{Merin's relevance and Carnap's relevance theory}
\label{sec:p15-w02b-p15-024-carnap}

Merin credits two of his definitions to Carnap (1950): extreme
positivity, ``modified'', and the covariance measure \(R_C\)
\cite[pp.~266, 272]{Merin2021Relevance}. Carnap's own theory,
reconstructed in Section~\ref{sec:p15-carnap-relevance} and
Section~\ref{sec:p15-carnap-relevance-analysis}, differs from Merin's
in four respects that the attributions do not show.
\begin{enumerate}
\item \emph{Prior evidence.} Carnap's relevance is three-place: an
  observation \(i\) is relevant to \(h\) \emph{on evidence} \(e\), and
  his relevance measure \(r(i,h,e)\) carries \(e\) as an argument.
  Merin's relations and \(R_C\) are two-place relative to \(P\); they
  correspond to Carnap's initial relevance, on tautological evidence,
  or to Carnap's relevance on \(e\) after conditioning on \(e\).
\item \emph{Contradictory evidence.} Carnap's definition D65-1d counts
  \(i\) irrelevant when \(e\land i\) is logically false, a condition on
  sentences. Merin's Definition~2 counts \(A\) irrelevant when
  \(P(A)=0\), a condition on the measure. Carnap remarks that in
  infinite systems his clause cannot be replaced by the condition
  \(m(e\land i)=0\), because an almost L-false observation can have
  zero relevance measure and still raise a defined conditional
  confirmation \cite[\S65, pp.~347--348]{Carnap1950Probability}
  \cite[p.~266]{Merin2021Relevance}.
\item \emph{Extreme positivity.} Carnap defines it structurally: \(i\)
  is positive and no stronger sentence is negative (D74-1). That the
  addition of \(i\) raises the confirmation of \(h\) to one, from below
  one, is a theorem for finite systems (T74-1(i)). Merin takes the
  probability-one condition as the definition and applies it to any
  probability, regular or not \cite[\S74, pp.~410--412]{Carnap1950Probability}.
\item \emph{Names and evaluation.} The relevance quotient Q is Carnap's
  name (note~8), but Carnap reports the quotient only as Johnson's and
  Keynes's tool and builds his theory on \(r\), praising its two
  additivity laws. Merin's verdict on the same two laws is the
  opposite: the conjunction law is neither useful nor intuitive, and the
  disjunction law is very counterintuitive, because a disjunction's
  relevance should lie between that of its disjuncts
  \cite[\S67, pp.~361--362]{Carnap1950Probability}
  \cite[pp.~272--273]{Merin2021Relevance}.
\end{enumerate}

\begin{hbtheorem}{thm:p15-w02b-p15-024-carnap-bridge}{proved=annex:p15-w02b-p15-024-carnap-bridge}{Carnap's relevance measure and Merin's covariance measure}
Let \(m\) be a regular measure function in Carnap's sense on a finite
system, \(c(h,x)=m(hx)/m(x)\) its confirmation function, and \(e\) a
sentence with \(m(e)>0\). Write \(P_e=m(\cdot\land e)/m(e)\).
\begin{enumerate}
\item \(r(i,h,e)=m(e)^2\,\Delta_{P_e}(i,h)\). For tautological \(e\),
  \(r(i,h,e)=R_C(i\Yright h)_m\).
\item \(i\) is positive, negative or irrelevant to \(h\) on \(e\) in the
  sense of D65-1 exactly when \(i\) is respectively
  \(\mathrm{pos}_{P_e}\), \(\mathrm{neg}_{P_e}\) or \(\perp_{P_e}\)
  to \(h\). Carnap's two additivity laws are Merin's \(R_C\) laws of
  Theorem~\ref{thm:p15-r43-composition} for the measure \(P_e\).
\item For tautological \(e\), \(i\) is extremely positive to \(h\) in
  the sense of D74-1 exactly when \(i\,\mathrm{expos}_m\,h\). For an
  arbitrary probability on a finite set, \(A\,\mathrm{expos}_P\,B\)
  holds exactly when \(P(A\neg B)=0\), \(P(AB)>0\) and
  \(P(\neg A\neg B)>0\). If \(P\) is not regular this can hold although
  \(A\) does not entail \(B\): on \(\{1,2,3\}\) with masses
  \((1/2,0,1/2)\), \(A=\{1,2\}\) is extremely positive to
  \(B=\{1\}\).
\item \(R_C\) satisfies MDHD whenever \(P(E)>0\), and violates POP-D.
\end{enumerate}
\end{hbtheorem}

Item~4 places Carnap's measure among the axioms. Crupi and Tentori
note that it satisfies their disjunction axiom (C2) while departing
from (C1) \cite[p.~656, n.~9]{CrupiTentori2016Confirmation}. In Merin's
terms it shares MDHD with D but lacks the axiom he treats as the entry
condition for a first-resort measure. Because
\(R_C(E\Yright H)_P=P(E)\,R_D(E\Yright H)_P\), it agrees in sign with D
and disagrees in ranking whenever the sources have different
probabilities.

\subsection{The raven ordering and Hempel's objection}
\label{sec:p15-w02b-p15-024-raven}

Merin's raven solution, in Section~\ref{sec:p15-r43-rankings}, uses
Hosiasson-Lindenbaum's assumptions: the hypothesis is independent of
being a raven and of being nonblack, and there are far fewer ravens
than nonblack things \cite[pp.~271--272]{Merin2021Relevance}. Hempel,
in the footnote discussed in
Section~\ref{sec:p15-hempel-raven-responses}, accepted that from her
axioms the smaller increment for a nonblack nonraven follows
\emph{provided} ravens are few compared with nonblack objects, but
denied that this proviso is warranted in general. In a co-ordinate
language whose individuals are finite space-time regions, the regions
containing a raven and those containing something nonblack both have
the power of the continuum; in a language of physical things the
assumption may well hold \cite[I, p.~21, n.~2]{Hempel1945ConfirmationI}.
Merin does not discuss this objection. The following family shows that
it bears directly on the ordinal conclusion.

\begin{hbtheorem}{thm:p15-w02b-p15-024-raven-imbalance}{proved=annex:p15-w02b-p15-024-raven-imbalance}{The raven ordering is equivalent to the imbalance assumption}
Fix \(P(H)\in(0,1)\) and numbers \(\rho,\nu,\varepsilon\) with
\(\rho+\nu\leq1\) and \(0<\varepsilon<\min\{\rho,\nu\}\). Given
\(H\), let a sampled object be a raven with probability \(\rho\) and
nonblack with probability \(\nu\), never both. Given \(\neg H\), keep
these two marginals and let both occur with probability \(\varepsilon\).
Let \(A\) say the object is a black raven and \(B\) that it is a
nonblack nonraven. Then both are positive to \(H\), and
\[
 R_L(A\Yright H)_P=\frac{\rho}{\rho-\varepsilon},\qquad
 R_L(B\Yright H)_P=\frac{\nu}{\nu-\varepsilon}.
\]
Hence \(A\) is more relevant to \(H\) than \(B\), equally relevant,
or less relevant, on every measure satisfying POP-D, according as
\(\rho<\nu\), \(\rho=\nu\) or \(\rho>\nu\).
\end{hbtheorem}

The model of Section~\ref{sec:p15-r43-rankings} is the case
\(\rho=1/100\), \(\nu=4/5\), \(\varepsilon=1/200\). Theorem~9 of
Merin's chapter (here Theorem~\ref{thm:p15-r43-rankings}) makes the
ordering measure-independent once the model is fixed, but the model
itself has to supply \(\rho<\nu\). Where Hempel's objection applies and
no such imbalance is available, the ordering has no basis in this
family. Hempel's own satisfaction criterion, a qualitative relation,
is a separate matter treated in Chapter~\ref{sec:p15-hempel-confirmation-and-paradoxes}.

\section{Old evidence, common causes and measure pluralism}
\label{sec:p15-w02b-p15-024-cited-routes}

Three passages of Merin's chapter summarize positions of other writers
in a sentence each: the routes to the problem of old evidence
(p.~268), the common-cause reading of screening and Pearl's thesis
about diagnostic categories (p.~267), and the recommendation of measure
pluralism (p.~274). None of the works concerned is reproduced here, so
their arguments are not reconstructed. What can be done, and is done
below, is to state exactly what Merin attributes, to prove the
probability facts on which each summary turns, and to show where those
facts limit the summarized position.

\subsection{The routes to the problem of old evidence}
\label{sec:p15-w02b-p15-024-old-evidence}

Merin's setting is a hypothesis \(H\) of positive probability that
entails a possible data report \(D\). If \(D\) is already known and
\(P\) is required to record that knowledge, then \(P(D)=1\) and \(D\)
is irrelevant to \(H\); Merin attaches this predicament to Glymour
(1980) and illustrates it with general relativity and the long-known
anomaly in the orbit of Mercury, which the theory entailed and which
was widely judged to raise confidence in it. He lists three ways in
which the judgment has been reconstructed and a fourth way out
\cite[p.~268]{Merin2021Relevance}:
\begin{enumerate}
\item update histories that are hypothetical and in part not obtained
  by conditioning, with a reference to Jeffrey (2004);
\item probabilities over primitive propositions saying that \(H\)
  entails or supports \(D\), with references to Earman (1992) and to
  Fitelson and Hartmann (2015);
\item leaving \(D\) slightly uncertain and choosing a relevance measure
  that still gives near-certain \(D\) an appreciable amount of
  relevance, with a reference to H\'ajek and Joyce (2008);
\item the statisticians' practice of replacing the credal requirement
  by a model of conditional frequencies for hypothetical repeatable
  observations, for which Merin names no single work; note~5 points to
  Edwards (1972) on statistical likelihood.
\end{enumerate}
The irrelevance of known evidence is proved in
Theorem~\ref{thm:p15-r43-polarity}, and the fourth route is the subject
of Theorem~\ref{thm:p15-w02b-p15-024-certain-data}. The third route
names no measure. The next result shows that it cannot be any of the
measures defined in Chapter~\ref{sec:p15-rationality-4-3} read in the
direction from \(D\) to \(H\).

\begin{hbtheorem}{thm:p15-w02b-p15-024-near-certain}{proved=annex:p15-w02b-p15-024-near-certain}{Near-certain entailed evidence on the principal measures}
Let \(H\) entail \(D\), and let \(0<P(H)<P(D)<1\). Write \(h=P(H)\)
and \(d=P(D)\). Then \(D\) is positive to \(H\), and
\[
 \begin{aligned}
 R_D(D\Yright H)_P&=\frac{h(1-d)}{d},&
 R_Q(D\Yright H)_P-1&=\frac{1-d}{d},\\
 R_L(D\Yright H)_P-1&=\frac{1-d}{d-h},&
 R_Z(D\Yright H)_P&=\frac{h(1-d)}{d(1-h)}.
 \end{aligned}
\]
For fixed \(h\), each of these tends to the measure's neutral value as
\(d\to1\), and so do \(R_{QL}\), \(R_{LL}\) and \(R_{KO}\). By
contrast, \(R_Z(H\Yright D)_P=1\), the largest value of Z, for every
such \(P\), and \(R_L(\neg D\Yright H)_P=0\), so that
\(R_{LL}(\neg D\Yright H)_P=-\infty\).
\end{hbtheorem}

On D, Q, L and Z the relevance of the entailed evidence to the
hypothesis is proportional to \(1-d\) up to a factor bounded on
\(d\geq(1+h)/2\), so keeping \(D\) only slightly uncertain makes it
only slightly relevant. An appreciable value for near-certain \(D\)
appears in this family only when the direction is reversed, from \(H\)
to \(D\) on Z, or when the evidence is negated. Whether H\'ajek and
Joyce propose one of these readings, or a measure outside the family,
cannot be settled from Merin's sentence, and their text is not
reproduced here. The theorem fixes what any such proposal has to
supply: a measure that is not a function of the source--target joint
table vanishing with \(1-d\), or an explicit argument for reading the
relevance from hypothesis to evidence.

The first two routes change the model rather than the measure. A
history that is not obtained by conditioning, or an algebra enlarged
by propositions about entailment, is a different input from the single
credal state \(P\) with \(P(D)=1\); the results of this chapter and of
Chapter~\ref{sec:p15-rationality-4-3} apply to such a reconstruction
only after its probability model has been written down. Merin gives
none, and none is supplied here. The statisticians' route is the one
on which Merin comments (``an easier life''), and the comment is
qualified by Theorem~\ref{thm:p15-w02b-p15-024-certain-data}: the
discriminating likelihood is a property of the frequency model, not of
the scientist's credal state.

\subsection{Common causes and diagnostic categories}
\label{sec:p15-w02b-p15-024-common-cause}

Merin states his Theorem~6, that under conditional independence given
\(B\) and given \(\neg B\) mutual positivity of \(A\) with \(B\) and of
\(B\) with \(C\) makes \(B\) and the conjunction \(AC\) mutually
positive, as a special case of a result shown by Reichenbach (1956).
Note~3 states that result as a product rule for polarities and adds
that the screening premise can be weakened to absence of conditional
negativity, crediting the weakening to W.~E. Johnson by way of Keynes
(1921) \cite[p.~267 and p.~274, n.~3]{Merin2021Relevance}. The
probability content of both is proved in
Theorem~\ref{thm:p15-r43-screening}. Reichenbach's intended reading,
with \(B\) a common cause of \(A\) and \(C\), is an interpretation of
that theorem: the theorem holds for every \(B\) that screens, causal or
not, and a causal reading needs the separate premises set out in
Chapter~\ref{sec:p16-finite-causal-models}. Merin also notes that
without the screening premise a conjunction of two positive sources can
be irrelevant or negative to the target, citing Carnap (1950,
chapter~6). Carnap's complete sign classification,
Theorem~\ref{thm:p15-carnap-relevance-sign-cases}, shows that two
positive observations can have a conjunction of either other sign, and
item~2 of Theorem~\ref{thm:p15-w02b-p15-024-carnap-bridge} carries his
signs over to Merin's relations.

Merin reports Pearl's (1988) surmise that medical diagnostic categories
are formed so that symptoms become conditionally independent given the
category and given its absence, the category being identified as their
common cause. He proposes a test: search for diseases whose
symptomatology is paradoxical in Carnap's sense, and he argues that
credence in Pearl's thesis should rise if none is found, because
finding one would be extremely negative to the thesis
\cite[p.~267]{Merin2021Relevance}. The thesis is a claim about how a
body of medical classification came to be, and Merin gives no data for
or against it. His argument for the test is formal and can be made
exact.

\begin{hbtheorem}{thm:p15-w02b-p15-024-search}{proved=annex:p15-w02b-p15-024-search}{What an unsuccessful search for counterexamples shows}
Let \(T\) be a thesis and \(F\) the report that a search finds a
counterexample, with \(0<P(T)<1\), \(0<P(F)<1\) and \(P(TF)=0\). Then
\(F\,\mathrm{exneg}_P\,T\), \(\neg F\) is positive to \(T\), and
\[
 R_D(\neg F\Yright T)_P=\frac{P(T)\,P(F)}{1-P(F)},\qquad
 R_Q(\neg F\Yright T)_P=\frac{1}{1-P(F)},
\]
\[
 R_L(\neg F\Yright T)_P=
 \begin{cases}
 1/\bigl(1-P(F\mid\neg T)\bigr),&P(F\mid\neg T)<1,\\
 +\infty,&P(F\mid\neg T)=1.
 \end{cases}
\]
\end{hbtheorem}

The sign of the conclusion is Merin's: a failed search speaks for the
thesis. Its size is fixed by \(P(F\mid\neg T)\), the probability that
the search would find a counterexample if the thesis were false. A
search with little chance of success leaves \(R_L\) near one and
\(R_D\) near zero, however long it runs. The premise \(P(TF)=0\) also
matters. If the thesis is read as a tendency, so that one paradoxical
disease does not refute it, \(F\) is at most negative to \(T\); then
\(\neg F\) is still positive to \(T\) by the polarity laws of
Theorem~\ref{thm:p15-r43-polarity}, but the three formulas no longer
hold and the size of the increment depends on the joint model of
thesis and search.

Note~4, attached to Pearl's thesis, names what conditional independence
of symptoms given a category buys: the specific benefit is
relevance-additivity under Lemma~3 \cite[p.~267 and p.~274,
n.~4]{Merin2021Relevance}. In the thesis the category \(B\) is the
target and the symptoms \(A\) and \(C\) are the sources, so Lemma~3
applies with \(H=B\), \(E_1=A\) and \(E_2=C\). The result below states
the link and shows that the same premise is what excludes the
paradoxical symptomatologies the test searches for.

\begin{hbtheorem}{thm:p15-w02b-p15-024-diagnostic-additivity}{proved=annex:p15-w02b-p15-024-diagnostic-additivity}{Symptoms screened by a category combine additively}
Let \(0<P(B)<1\) and \((A\perp C\mid\pm B)_P\), and write
\(a_1=P(A\mid B)\), \(b_1=P(A\mid\neg B)\), \(a_2=P(C\mid B)\),
\(b_2=P(C\mid\neg B)\).
\begin{enumerate}
\item If \(a_1,a_2,b_1,b_2>0\), then
  \[
   \begin{aligned}
   R_L(AC\Yright B)_P&=R_L(A\Yright B)_P\,R_L(C\Yright B)_P,\\
   R_{LL}(AC\Yright B)_P&=R_{LL}(A\Yright B)_P+R_{LL}(C\Yright B)_P .
   \end{aligned}
  \]
\item If \(A\) and \(C\) are each positive to \(B\), then \(AC\) is
  positive to \(B\). If moreover \(b_1,b_2>0\), then
  \(R_{LL}(AC\Yright B)_P>\max\{R_{LL}(A\Yright B)_P,R_{LL}(C\Yright B)_P\}>0\).
\item Hence a symptom pair in Carnap's paradoxical configuration, with
  \(A\) and \(C\) each positive to \(B\) and \(AC\) irrelevant or
  negative to \(B\), refutes \((A\perp C\mid\pm B)_P\).
\end{enumerate}
\end{hbtheorem}

Item~2 is the probability content of Merin's Theorem~6 for this
configuration (Theorem~\ref{thm:p15-r43-screening}), here obtained from
the additive form: on the LL scale each screened positive symptom adds
a positive amount, so the conjunction cannot fall to zero or below.
Item~3 is why a paradoxical symptomatology counts against the thesis:
it is incompatible with the screening that the thesis says a diagnostic
category is formed to secure. It refutes the thesis for that category
only if the thesis is read as requiring exact screening; under a
tendency reading the remarks after
Theorem~\ref{thm:p15-w02b-p15-024-search} apply. Item~1 is the benefit
note~4 names: with the odds form of Bayes' theorem
(Theorem~\ref{thm:p15-r43-measures}), the posterior log-odds of the
category given both symptoms is its prior log-odds plus the two LL
values, each computed from that symptom alone. Without the screening
premise the conjunction's LL value must be obtained from the joint
likelihoods of the symptom pair.

\subsection{Disjoint disjuncts and mutual extreme negativity}
\label{sec:p15-w02b-p15-024-disjuncts}

Merin's Theorem~10 assumes \(P(AB)=0\) and shows that the relevance of
\(A\vee B\) to a target is a weighted mean of the relevances of \(A\)
and \(B\) on D and Q, and on L when neither disjunct is extremely
relevant to the target (Theorem~\ref{thm:p15-r43-composition}).
His conclusion describes the condition differently, as mutual extreme
negativity of the disjuncts \cite[pp.~273--274]{Merin2021Relevance}.
The two descriptions agree, with one degenerate exception in which the
mixture is trivial.

\begin{hbtheorem}{thm:p15-w02b-p15-024-exneg-disjuncts}{proved=annex:p15-w02b-p15-024-exneg-disjuncts}{Disjointness and mutual extreme negativity}
Let \(A\), \(B\) and \(C\) be events under \(P\).
\begin{enumerate}
\item If \(A\,\mathrm{exneg}_P\,B\), then \(P(A)>0\), \(P(B)>0\),
  \(P(AB)=0\) and \(B\,\mathrm{exneg}_P\,A\). Extreme negativity is
  therefore symmetric.
\item If \(P(A)>0\) and \(P(B)>0\), then \(P(AB)=0\) holds if and only
  if \(A\,\mathrm{exneg}_P\,B\), and so if and only if \(A\) and \(B\)
  are mutually extremely negative.
\item If \(P(AB)=0\) and \(P(A)=0<P(B)\), then for \(T=D,Q,L\)
  \(R_T(A\vee B\Yright C)_P=R_T(B\Yright C)_P\), and the mixture weight
  of Theorem~10 is \(\alpha=0\) for D and Q, and also for L whenever
  \(0<P(C)<1\) and \(P(B\mid\neg C)>0\).
\end{enumerate}
\end{hbtheorem}

So outside the case in which one disjunct has probability zero, and
the disjunction carries exactly the relevance of the other disjunct,
the premise of Theorem~10 is the premise that each disjunct
conclusively excludes the other. The first item also shows that the
word ``mutual'' adds nothing to the relation: one direction implies the
other. The case Merin calls very counterintuitive, Carnap's covariance
measure adding the relevances of disjoint disjuncts, rests on the same
premise; the mixture property is what distinguishes the three principal
measures and their increasing transforms from it.

\subsection{Measure pluralism}
\label{sec:p15-w02b-p15-024-pluralism}

Merin's conclusion recommends LL as the default when relevance must be
quantified, and Q or QL where direction is intrinsically unspecified.
It then records two positions that pull the other way, and refers to
H\'ajek and Joyce (2008) for both: that measures violating POP-D have
been commended for special purposes, and that some philosophers
recommend a plurality of measures to express different aspects of
relevance \cite[p.~274]{Merin2021Relevance}. Neither the special
purposes nor the commended measures are named. The formal results of
these two chapters bound what pluralism can change. Among measures
satisfying POP-D, rankings of several sources for one target do not
depend on the measure (Theorem~\ref{thm:p15-r43-rankings}) or on the
direction of measurement
(Theorem~\ref{thm:p15-w02b-p15-024-direction}), so pluralism within
that class makes a difference only to comparisons of targets, to
amounts, and to the composition rules. The measures that violate POP-D
include Carnap's \(R_C\) and the two likelihood differences of
Theorem~\ref{thm:p15-w02b-p15-024-z-characterization}; admitting them
makes source rankings depend on the probability of the source.
Whether that dependence serves a special purpose is a normative
question that the cited work, not reproduced here, would have to
answer.

\subsection{The surveys and the psychological pointers}
\label{sec:p15-w02b-p15-024-surveys}

Note~18 attaches to the remark that relevance has uses in philosophy of
science and lists surveys of them: Huber (2007), H\'ajek and Joyce
(2008), Crupi (2015) and Crupi and Tentori (2016)
\cite[p.~273 and p.~274, n.~18]{Merin2021Relevance}. The note makes no
claim beyond the pointer. Of the four, the survey by Crupi and Tentori
is used throughout this chapter at the pages cited; the other three are
not reproduced, and nothing here depends on them.

Note~19 attaches to the claim that, in some contexts, participants
asked for posterior probabilities or odds give answers consistent with
covert relevance judgments, the prominent instances being base-rate
neglect and the conjunction fallacy. It refers to Kahneman (2011,
chapters~15--16) on the two phenomena and to the chapters by Evans and
by Hertwig and Kozyreva in the same handbook
\cite[p.~273 and p.~274, n.~19]{Merin2021Relevance}; those two
chapters are presented in Chapters~\ref{sec:p13-rationality-1-2}
and~\ref{sec:p17-rationality-8-5}. Merin's text reports no task,
population, design, measure or uncertainty for either phenomenon, and
Kahneman's book is not reproduced here. The claim is therefore a
reading that the arithmetic permits, set out in
Section~\ref{sec:p15-r43-use}, not a finding; the original task
designs and their results are presented in
Chapters~\ref{sec:p19-conjunction-judgments},
\ref{sec:p19-rationality-2-4}
and~\ref{sec:p19-natural-frequencies-bayesian-reasoning}.

\section{The relative-distance measure and the dependent results}
\label{sec:p15-w02b-p15-024-dependents}

Merin's note~13 directs the reader to Crupi and Tentori (2016) for
measures other than D, Q and L that satisfy POP-D, Z among them, and to
Crupi and Tentori (2013) for a uniqueness result for Z
\cite[p.~274, n.~13]{Merin2021Relevance}. The 2016 survey states that
result as its Theorem~2 and refers to the 2013 paper for the proof. Its
two further axioms, written in Merin's notation, are as follows
\cite[pp.~658--659]{CrupiTentori2016Confirmation}.
\begin{description}
\item[Partial inconsistency (C5).] If \(P(EH)\leq P(E)P(H)\), then
  \(R(E\Yright H)_P=R(H\Yright E)_P\).
\item[Confirmation complementarity (C6).]
  \(R(E\Yright H_1)_P\bowtie R(E\Yright H_2)_P\) exactly when
  \(R(E\Yright\neg H_2)_P\bowtie R(E\Yright\neg H_1)_P\), where
  \(\bowtie\) is any of \(<,=,>\).
\end{description}
The first makes relevance symmetric where source and target are
negatively correlated, as logical inconsistency is symmetric. The second
ties support for a hypothesis inversely to the support for its negation.
Crupi and Tentori relate it to Keynes's remark that an argument is as
near to proving a proposition as it is to disproving its contradictory
\cite[p.~659]{CrupiTentori2016Confirmation}.

\begin{hbtheorem}{thm:p15-w02b-p15-024-z-characterization}{proved=annex:p15-w02b-p15-024-z-characterization}{Characterization of the relative-distance family}
Let \(R\) be given by one formal function of the joint table, uniformly
over all finite probability models, and compared only on events with
probabilities strictly between zero and one.
\begin{enumerate}
\item \(R\) satisfies POP-D, partial inconsistency and confirmation
  complementarity if and only if \(R=f\circ R_Z\) for a strictly
  increasing \(f\) on \([-1,1]\).
\item \(R_Z(E\Yright\neg H)_P=-R_Z(E\Yright H)_P\).
\item Each of \(P(E\mid H)-P(E)\), \(P(E\mid H)-P(E\mid\neg H)\) and
  \(R_C(E\Yright H)_P\) equals \(P(E)\) times a function of \(P(H)\)
  and \(P(H\mid E)\), is positive exactly when \(R_D\) is, and
  violates POP-D.
\end{enumerate}
\end{hbtheorem}

Item~3 covers the three measures Crupi and Tentori set aside as
departing from their first axiom, those of Mortimer, Nozick and Carnap
\cite[p.~656, n.~9]{CrupiTentori2016Confirmation}. Among the
incremental measures the survey discusses, those satisfying POP-D are
the increasing transforms of D, of Q (the survey lists five variants),
of L, and of Z \cite[pp.~655, 658]{CrupiTentori2016Confirmation}.
The posterior probability \(P(H\mid E)\) also satisfies POP-D, trivially,
but it measures how credible \(H\) is after \(E\), not how much \(E\)
changed that credibility. It has no neutral value separating support
from countersupport: the posterior \(1/2\) is positive relevance when
\(P(H)=1/4\) and negative relevance when \(P(H)=3/4\). It therefore
fails Merin's Quality Representation, and
Crupi and Tentori note that the condition characterizing it is
inconsistent with their discrimination requirement
\cite[p.~269]{Merin2021Relevance}
\cite[p.~656, n.~10]{CrupiTentori2016Confirmation}.

\begin{proofboundary}{def:p15-w02b-p15-024-dependents}{Standing of the results Merin's chapter relies on}
Every theorem and lemma of Merin's chapter other than Lemma~4 is
proved in the annex to Chapter~\ref{sec:p15-rationality-4-3}, including the realization
argument for his Theorem~7 that follows Crupi, Chater and Tentori (2013)
and Crupi and Tentori (2016). The directions of Theorem~7 that Merin's
note~12 credits to Steel (2003), that each displayed family satisfies
its axioms, are the converse clauses of that proof. The uniqueness
result for Z of note~13 is
Theorem~\ref{thm:p15-w02b-p15-024-z-characterization}, proved in this
chapter's annex without the 2013 paper. The two inferences of Merin's
axiom-status argument are
Theorem~\ref{thm:p15-w02b-p15-024-axiom-status}, and Lemma~4, whose
content appears there only inside Theorem~\ref{thm:p15-r43-rankings},
is stated precisely and proved as
Theorem~\ref{thm:p15-w02b-p15-024-direction}.

The works on which the chapter's claims depend form a finite set: the
thirty-three entries of its reference list, together with seven chapters
of the same handbook named in its text and notes
\cite[pp.~267, 273 and pp.~274--275]{Merin2021Relevance}. Each has one
of the following standings.
\begin{description}
\item[Formal content proved here or in the annex to
  Chapter~\ref{sec:p15-rationality-4-3}.] Crupi, Chater and Tentori
  (2013) and Crupi and Tentori (2016), for Theorem~7; Crupi and Tentori
  (2013), for the uniqueness of Z, proved independently; Steel (2003),
  for the converse clauses of note~12; Reichenbach (1956), for the
  product rule of note~3 and Theorems~5 and~6, proved in
  Theorem~\ref{thm:p15-r43-screening}; Carnap (1950), for Definition~4,
  the measure \(R_C\), compared with his text in
  Section~\ref{sec:p15-w02b-p15-024-carnap}, and the paradoxical
  conjunctions of his chapter~6, given by
  Theorem~\ref{thm:p15-carnap-relevance-sign-cases}; Hosiasson-Lindenbaum
  (1940, p.~138), for the raven assumptions, whose consequences are
  Theorem~\ref{thm:p15-w02b-p15-024-raven-imbalance}; Good (1960) and
  Heckerman (1988), for MCIE and Lemma~3; Savage (1954, pp.~128--140),
  for the sufficiency of L, proved in the finite binary case; Merin
  (1999), for the proposition of note~16.
\item[Arguments attributed but not given.] Steel (2003), that POP-D is
  necessary for enumerative induction; Festa and Cevolani (2016), the
  Matthew properties of note~17; Merin (1999), the function-word
  semantics of note~20; Glymour (1980), the statement of the old-evidence
  problem, whose probability core is proved; Jeffrey (2004), Earman
  (1992) and Fitelson and Hartmann (2015), the first two old-evidence
  routes; H\'ajek and Joyce (2008), the third route, the commendation of
  measures violating POP-D and measure pluralism, bounded by
  Theorem~\ref{thm:p15-w02b-p15-024-near-certain} and
  Section~\ref{sec:p15-w02b-p15-024-pluralism}; Pearl (1988), the thesis
  on diagnostic categories, whose proposed test is
  Theorem~\ref{thm:p15-w02b-p15-024-search}; Edwards (1972), statistical
  likelihood (note~5); Eells and Fitelson (2002), the number of measures
  and the prevailing view on negation (pp.~269, 272). None of these
  arguments is used in a proof.
\item[Pointers to surveys.] Huber (2007), H\'ajek and Joyce (2008),
  Crupi (2015) and Crupi and Tentori (2016), listed in note~18; Earman
  (1992), H\'ajek and Joyce (2008) and Crupi (2015) are also cited as
  accessible accounts of the raven solution (p.~271). Only the last
  survey is used here.
\item[Historical attributions.] Keynes (1921), for MacColl's \(\delta\),
  Johnson's coefficient of influence and Johnson's weakening in note~3;
  Keynes (1921), Carnap (1950) and Jeffrey (2004), as users of the word
  \emph{relevance} (p.~268); Good (1950), Kemeny and Oppenheim (1952)
  and Good (1983), for QL, KO and the term quasi-utility; Savage (1954),
  for the comparison with expected utility (p.~273). These rest on the
  attributions given in Section~\ref{sec:p15-w02b-p15-024-history}.
\item[Empirical pointers.] Tentori, Crupi and Russo (2013) and Cevolani
  and Crupi (2015) on conjunction judgments; Wilson and Sperber (2004)
  and Sperber, Cara and Girotto (1995) on relevance theory and the
  selection task; Skovgaard-Olsen, Singmann and Klauer (2016a, 2016b) on
  relevance and conditionals \cite[p.~273]{Merin2021Relevance};
  Kahneman (2011, chapters~15--16) in note~19. The chapter reports no
  task, population, design, measure or uncertainty for any of them, so
  they support no conclusion here; their separation from the relevance
  mathematics is set out in Section~\ref{sec:p15-r43-use} and
  Section~\ref{sec:p15-w02b-p15-024-surveys}.
\item[Chapters of the same handbook.] H\'ajek and Staffel (4.1), Dubois
  and Prade (4.7) and Kern-Isberner, Skovgaard-Olsen and Spohn (5.3),
  named in note~1, and Evans (1.2) and Hertwig and Kozyreva (8.5), named
  in note~19, are presented in the chapters linked in
  Sections~\ref{sec:p15-w02b-p15-024-notes}
  and~\ref{sec:p15-w02b-p15-024-surveys}. Pearl's chapter (7.1), cited
  for Simpson's paradox and for the diagnostic thesis (p.~267), and
  Peterson's chapter (8.2), cited for expected utility (p.~273), have no
  presentation in this handbook; nothing here depends on them.
\end{description}
\end{proofboundary}

\index{evidential relevance!axioms}\index{relevance quotient}
\index{relevance measure (Carnap)}\index{relative distance measure}
\index{raven paradox!Hempel's objection}\index{likelihood!law of}
